\section{Conclusion}
\label{sec:conclusion}

This study shows that discrimination alone is not sufficient to establish the reliability of autism classification from resting-state functional connectomes. Using a fixed analytical framework, the effects of cohort composition, motion, regularisation, probability calibration, and model interpretation are examined alongside predictive performance. The ensemble achieves an AUROC of 0.776 and maintains similar discrimination when acquisition sites are withheld, but its performance varies with the cognitive composition of the autistic sample. Classifier scores separate diagnostic groups without showing a corrected association with symptom severity within autistic participants. Lower somatomotor connectivity also persists and strengthens after within-site matching on recorded mean FD, although this remains an exploratory group-level finding rather than evidence of a biomarker. At the modelling level, regularised linear classifiers perform better than the evaluated fixed-configuration tree ensembles, while the apparent advantage of windowed connectivity disappears when covariance regularisation is matched. The native ensemble probabilities show a systematic calibration offset that is corrected using training-only recalibration without reducing discrimination. Interpretation is also more reliable at the level of functional systems: default-mode pairs contribute the largest overall share of classifier separation, while system-level attribution reproduces more consistently than individual connection patterns. Taken together, these findings support a view of trustworthy connectome-based autism classification in which predictive performance is interpreted together with calibration, cohort characteristics, analytical design, and the reproducibility of model explanations. Independent external validation is required before these conclusions can be extended beyond the present retrospective research cohort.